% TeX source
%
% Author: Tetsuya Ishikawa <tiskw111@gmail.com>
% Date  : Aug  4, 2019
%%%%%%%%%%%%%%%%%%%%%%%%%%%%%%%%%%%% SOURCE START %%%%%%%%%%%%%%%%%%%%%%%%%%%%%%%%%%%

\documentclass[uplatex, twocolumn, a4paper, 10pt, dvipdfmx]{jsarticle}

\usepackage[dvipdfm,width=190truemm,height=270truemm]{geometry}
\usepackage{amsmath,amssymb}
\usepackage{lmodern}
\usepackage{helvet}
\usepackage[sourcehan, unicode]{pxchfon}
\usepackage[dvipdfmx]{graphicx}
\usepackage{thmbox}
\usepackage{csquotes}

\allowdisplaybreaks

\newcommand{\bs}[1]{\boldsymbol{#1}}
\newcommand{\mbar}{\hspace{1pt}\bigm|\hspace{1pt}}
\newcommand{\tran}{^{\mkern-1.5mu\mathsf{T}}}
\DeclareMathOperator*{\sgn}{sgn}
\DeclareMathOperator*{\Sgn}{Sgn}
\DeclareMathOperator*{\argmin}{argmin}

\newtheorem[M]{theorem}{定理}[section]
\newtheorem[M]{definition}[theorem]{定義}
\newtheorem[M]{lemma}[theorem]{補題}

\renewenvironment{proof}%
{\textsf{証明}：}{\hfill$\blacksquare$\vspace*{8pt}}

\newcommand{\titlebar}{
\setlength\unitlength{1truemm}
\vspace*{-19pt}
\begin{picture}(88,0)(0,0)
\put(-4,0){\line(1,0){93}}
\put(89,0){\circle*{1}}
\end{picture}}

\begin{document}

\title{\textsf{\bfseries Random Fourier Features}}
\author{石川 徹也 \quad \texttt{<tiskw111@gmail.com>}}
\date{2020年3月3日}
\maketitle

\section*{はじめに}
\titlebar
% TeX source
%
% Author: Tetsuya Ishikawa <tiskw111@gmail.com>
% Date  : October 13, 2021
%%%%%%%%%%%%%%%%%%%%%%%%%%%%%%%%%%%%%%%%%%%%%%%%%%%%% SOURCE START %%%%%%%%%%%%%%%%%%%%%%%%%%%%%%%%%%%%%%%%%%%%%%%%%%%%%

本文書では，ガウス過程モデル~\cite{Rahimi2007}に
random Fourier features~\cite{Rasmussen2006} を適用する手順を解説します．
これにより，ガウス過程モデルの学習や推論を高速化させることができ，
より大規模なデータにガウス過程モデルを適用することができます．

ガウス過程モデル~\cite{Rahimi2007}は確率的教師あり機械学習フレームワークのひとつであり，
サポートベクトルマシンやランダムフォレストなどと並んで，回帰や分類タスクに広く使用されています．
ガウス過程モデルがサポートベクトルマシンやランダムフォレストと大きく違う点は「確率的なモデルである」ことです．
すなわち，ガウス過程モデルは確率的なモデルとして定式化されているため，予測値だけでなく，
その予測に対する不確実性の尺度をも提供することができます．
これは機械学習の説明性を上げることのできる非常に有益な性質です．

その一方で，ガウス過程モデルは学習や推論の計算コストが高いことでも知られています．
学習データの総数を$N \in \mathbb{Z}^{+}$としたとき，
ガウス過程モデルの学習に要する計算量は$O(N^3)$，推論に要する計算量は$O(N^2)$です．
問題は計算量が学習データの総数$N$のべき乗になってしまっていることで，これは大規模データへの適用に際して障害になり得ます．
これはガウス過程モデルがカーネル法と同等の数学的構造を有していることに起因しており，
言い変えれば，カーネルサポートベクトルマシンも同じ悩みを有しています．

カーネル法を高速化する手法のひとつにrandom Fourier features~\cite{Rasmussen2006}があります（以下 RFF と略します）．
これはカーネル関数を有限次元ベクトルの内積として近似することで，
カーネル法の柔軟性を維持しつつ計算量を大幅に削減する手法です．
具体的には学習に要する計算量を$O(N D^2)$，推論に要する計算量を$O(D^2)$にまで削減することができます．
ただし$D \in \mathbb{Z}^{+}$は RFF のハイパーパラメータであり，
学習データの総数$N$とは独立に指定することができます．

ガウス過程モデルはカーネル法と同等の数学的構造を有しているため，ガウス過程モデルにも RFF を適用することができます．
これにより，ガウス過程モデルはより強力かつお手軽に使える，非常に頼もしいツールへと進化します．


しかしながら，ガウス過程モデルへのRFFの適用にあたっては，実は一筋縄ではいかないところがあります．
ただRFFを適用するだけでは高速化につながらず，一工夫する必要があるのです．
ですが，そのあたりの困難さや解決方法に言及している文献が，残念ながら世の中には存在しないようでしたので，
本文書にてその手順を解説しようと考えた次第です．

ちなみに，本文書は拙作ライブラリ\texttt{rfflearn}~\cite{rfflearn}に同梱されているドキュメントの日本語版です．
おそらく\texttt{rfflearn}に同梱されている文書の方が頻繁にメンテナンスされると思いますので，
英語でも差し支えない方はそちらをご参照下さい．

\begin{displayquote}\footnotesize\textsf{NOTE:}
    上述のライブラリ \texttt{rfflearn} は以下で公開されています．
    \begin{center}
        \texttt{https://github.com/tiskw/random-fourier-features}
    \end{center}
\end{displayquote}

%%%%%%%%%%%%%%%%%%%%%%%%%%%%%%%%%%%%%%%%%%%%%%%%%%%%% SOURCE FINISH %%%%%%%%%%%%%%%%%%%%%%%%%%%%%%%%%%%%%%%%%%%%%%%%%%%%
% vim: expandtab tabstop=4 shiftwidth=4 fdm=marker


\section{なぜRFFが必要なのか？}
\titlebar
% TeX source
%
% Author: Tetsuya Ishikawa <tiskw111@gmail.com>
% Date  : March 01, 2020
%%%%%%%%%%%%%%%%%%%%%%%%%%%%%%%%%%%%%%%%%%%%%%%%%%%%% SOURCE START %%%%%%%%%%%%%%%%%%%%%%%%%%%%%%%%%%%%%%%%%%%%%%%%%%%%%

機械学習におけるクラス分類問題を解く手法のひとつとして，カーネルサポートベクターマシンと呼ばれる手法がある．
これは通常のサポートベクターマシン（ここではこれを線形サポートベクターマシンと呼ぶことにしよう）にカーネル法を適用する
ことで柔軟性を格段に向上させたものであり，現在でも現役で使用されている，カーネル法の応用例のひとつでもある．
しかし，カーネルサポートベクターマシンには，学習や推論にかかる時間が学習データ数に多項式的に比例してしまうという決定的な弱点がある．
具体的には，学習データの数を$N$とすると，学習に要するステップ数は$O(N^3)$，推論に要するステップ数は$O(N)$である．
繰り返しになるが$N$は学習データの数である．これは，データの質よりも量の多さをよしとする機械学習の分野において，
致命的とも言える弱点である．この恐しさがお分かり頂けるであろうか？

想像して頂きたい．何らかのクラス分類器を作成したところ，性能があまり出なかったとしよう．
そうなったとき，学習データセットの拡充は間違いなく優先度の高い検討事項である．
しかし，先に述べた通り，カーネルサポートベクターマシンの場合は学習データの数を増やせば増やすほど
学習や推論にかかる時間が増えてしまう．
学習時間が増えるのはまだ許容範囲内であろうが，推論の時間まで増加してしまうのは許容し難いのではないだろうか．
もちろんケースバイケースではあるのだが，実応用では推論時間の上限は事前に決められており，また推論を行うデバイスも
変更できない場合が多く，推論の計算量の増加は致命的となるケースが多い．
つまり，カーネルサポートベクターマシンの場合は安易に学習データセットの追加ができないのである．

ちなみに，この弱点は線形サポートベクターマシンでは起こりにくい．
線形のサポートベクターマシンの場合，少なくとも推論時間は学習データセットの数によらず一定である．
つまり，この弱点はカーネル法によって高い柔軟性を得たことに由来するものなのである．
とは言え，線形サポートベクターマシンでは性能が出ないケースが多いのも事実であるから，
何とかしてカーネルサポートベクターマシンを高速化し，大量の学習データセットを適用したいものである．

この弱点を解消するために考え出されたのが，カーネル関数の有限次元近似であり，その中で最も成功しているのがRFFである．

これをキチンと説明するために，まずはカーネル法について簡単にご説明していこう．

%%%%%%%%%%%%%%%%%%%%%%%%%%%%%%%%%%%%%%%%%%%%%%%%%%%%% SOURCE FINISH %%%%%%%%%%%%%%%%%%%%%%%%%%%%%%%%%%%%%%%%%%%%%%%%%%%%
% vim: expandtab tabstop=4 shiftwidth=4 fdm=marker


\section{カーネルとは？}
\titlebar
% TeX source
%
% Author: Tetsuya Ishikawa <tiskw111@gmail.com>
% Date  : March 01, 2020
%%%%%%%%%%%%%%%%%%%%%%%%%%%%%%%%%%%%%%%%%%%%%%%%%%%%% SOURCE START %%%%%%%%%%%%%%%%%%%%%%%%%%%%%%%%%%%%%%%%%%%%%%%%%%%%%

本節ではカーネル法の概要を，数学的証明を用いずに説明してみたいと思う．
本節でカーネル法を説明する主な目的は，次節で行うRFFの解説において
事前に記号の理解を持って頂くことである．
厳密さが多少なりとも犠牲になるとについてはご容赦頂きたい．

やや深層学習に寄ったものの見方ではあるが，機械学習は主に
「データから特徴量を抽出する」ステップと「その特徴量を用いてタスクを行う」ステップに
大別される．図\ref{fig:kernel_method_and_feature_extraction}を参照されたい．
以下ではこの2つのステップを踏むことでカーネル法を概観してみよう．

\begin{figure}[t]
    \centerline{\includegraphics[width=200pt]{figures/kernel_method_and_feature_extraction.pdf}}
    \caption{カーネル法と特徴抽出}
    \label{fig:kernel_method_and_feature_extraction}
\end{figure}

\subsection{無限次元の特徴抽出と特徴関数}

入力データの集合を
$\mathcal{D} \triangleq \{ \bs{x}_i \in \mathbb{R}^{M} \bigm| i = 1, 2, \ldots, M \}$
とおく．
カーネル法は，入力データ$\bs{x}$から特徴量として関数，すなわち無限次元の特徴量を
抽出する手法である．このデータから無限次元の特徴量を抽出する関数，
ここれはそれを特徴関数と呼ぶこととするが，これを$\phi(\bs{x}, \bs{p})$で表すこととしよう．
具体例を挙げるとすれば，例えば
\begin{equation}
    \phi(\bs{x}, \bs{p}) \triangleq \exp \left( - \frac{\|\bs{x} - \bs{p}\|^2}{2 \sigma^2} \right),
    \label{eqn:feature_rbf}
\end{equation}
である．
これは有名なカーネル関数であるRFBカーネル (\textit{radial basis function kernel}) に対応する特徴関数である．
\begin{displayquote}\footnotesize\textsf{NOTE:}
後に説明するが，カーネル関数と特徴関数は\textgt{全くの別物}であるので注意されたい．
カーネル関数は特徴関数同士の内積として定義されるもので，特徴関数とは立場上全くの別物である．
ただ，RBFカーネルを含む有名どころのカーネル関数は結果として再生性をみたすゆえに
特徴関数とカーネル関数の関数形が一致し，説明する側としてはヤヤコシいのだが\ldots \verb|(^^;|
\end{displayquote}

注意して欲しいのは，特徴関数は入力データ$\bs{x}$を受け取り，$\bs{p}$に関する関数を返す関数である．
そういう意味で，特徴関数は$\phi(\bs{x}, \bs{p})$と書くよりも，$\phi_{\bs{x}}(\bs{p})$と書いた方が分かり良いであろう．
やや通常の関数表記から逸脱するが，以下では$\phi_{\bs{x}}(\bs{p})$の表記を用いることとする．
また，関数であることを特に強調したい場合は，引数の$\bs{p}$を省略して$\phi_{\bs{x}}$と書く．

\begin{displayquote}\footnotesize\textsf{NOTE:}
余談ではあるが，関数である$\phi_{\bs{x}}$を特徴「量」と呼称して良いものであろうか？
これはやや疑義を生ずる点であろうが，著者は量と呼称して問題ないと考えている．
関数は$L^2$空間に代表されるベクトル空間を構成する要素であるから，間違いなくベクトル「量」である．
\end{displayquote}

\subsection{特徴関数からのタスク実施}

さて，データ$\bs{x}$から特徴量$\phi_{\bs{x}}$が得られたところで，次はこの特徴量を用いて
回帰や分類などのタスクを行うことを考えていこう．
タスクの実施の仕方はタスクによって様々なのだが，最も一般的なアプローチは
\begin{equation}
    f(\bs{x}) = \idotsint_{\mathbb{R}^M} w(\bs{p}) \phi_{\bs{x}}(\bs{p}) \, 
    \mathrm{d}p_1 \cdots \mathrm{d}p_M,
    \label{eqn:task_kernel}
\end{equation}
という形で関数を構築する方法である．
回帰がしたいのであれば$f(\bs{x})$が目的の曲線を追従するように，
分類がしたいのであれば$f(\bs{x})$が対数尤度を出力するようにすれば良い．
ただし$w(\bs{p})$は学習パラメータに相当する関数である．
この関数を適切に学習させることで目的のタスクを実現するのがカーネル法である．

\begin{displayquote}\footnotesize\textsf{NOTE:}
式(\ref{eqn:task_kernel})が分かりづれぇ！と思ったそこの貴方．是非とも関数解析を習得されることをお薦めする．
式(\ref{eqn:task_kernel})はただの関数同士の「内積」であり，$f(\bs{x}) = \langle w, \phi_{\bs{x}} \rangle$と表記することもできる．
つまりやっていることは線形回帰と何も変わらず，ただ扱う空間がベクトル空間から関数空間に変わっただけのことである．
\end{displayquote}

もちろん，最適化パラメータが関数であるような最適化問題は数値計算問題として解きにくいので，
様々なテクニックを駆使してこれを有限次元の最適化問題に落し込み，その最適化を問題を解くことで学習を行うのである．

\begin{center}
メデタシ，メデタシ\ldots
\end{center}

\subsection{いやいや，カーネル関数はどこよ？}

え，カーネル関数はいつ出てくるのかって？
実はカーネル関数は，最適化問題をキチンと定式化し，その式をコネクリ回さないと出てこない．
本節ではその部分の理解にあまり体力を使いたくないので，先人が数式をコネクリ回した結果，以下が得られたということだけ指摘しておくに留めておく．
詳細については赤穂~\cite{Akaho2008}やBishop~\cite{Bishop2010}などを参照されたい．
\begin{enumerate}
    \item 関数$w$の最適化は有限次元のベクトルの最適化に落し込むことが出来る．ただしそのベクトルの次元は学習データ数$N$である．
          この関数の最適化問題を有限次元の最適化問題に落し込むテクニックは，カーネルトリックと呼ばれる．
    \item 上述1の最適化を解く上でも，最適化の結果を用いて推論を行う上でも，実は特徴関数$\phi_{\bs{x}}$は陽には出てこない．
          出てくるのは次の式
          \begin{equation}
              k(\bs{x}, \bs{y}) \triangleq \idotsint \phi_{\bs{x}}(\bs{p}) 
              \phi_{\bs{y}}(\bs{p}) \, \mathrm{d}p_1 \cdots \mathrm{d}p_M,
              \label{eqn:kernel}
          \end{equation}
          で与えられる関数$k$だけである．この関数$k$をカーネル関数と呼ぶ．
\end{enumerate}

さて，上記には2つの重要な事実が含まれている．
ひとつは，カーネルトリックによって関数次元の最適化問題を有限次元の最適化問題に落とし込めたとはいえ，
いまだに学習データ数$N$と同じだけの次元を持っているという点である．
これが本文書の冒頭で述べたカーネル法の弱点に直結しているのである．

もうひとつは，学習や推論を行う上で特徴関数は陽には出てこず，カーネル関数しか出てこないという点である．
ゆえに，一般的にはカーネル関数だけを設計し，特徴ベクトルには言及しない場合が多い．
ちなみに，式(\ref{eqn:feature_rbf})にて特徴関数の具体例をお見せしたが，これに対応するカーネル関数は
\begin{equation}
    k(\bs{x}, \bs{y}) = \exp(- \gamma \|\bs{x} - \bs{y}\|^2),
    \hspace{10pt} \gamma = - \frac{1}{4 \sigma^2},
    \label{eqn:kernel_rbf}
\end{equation}
である．これは式(\ref{eqn:kernel})を直接計算することで得られる．

さて，様々な式展開や証明を省略させて頂いたが，これでカーネル法のおさらいは終わりにして，早速RFFの定式化に移ろう．

\begin{displayquote}\footnotesize\textsf{NOTE:}
カーネル関数の説明が適当すぎないかって？
ご指摘の通りであるが，本文書としてはカーネル法の説明は前座にすぎないので，本節でパルプを3枚も4枚も費すことはしたくない．
もちろんカーネル関数がいきなり天から降ってくるのは本文書の想定読者を考えると論外である．
もっとスマートな導入の仕方があれば是非ご連絡頂きたい．
\end{displayquote}

%%%%%%%%%%%%%%%%%%%%%%%%%%%%%%%%%%%%%%%%%%%%%%%%%%%%% SOURCE FINISH %%%%%%%%%%%%%%%%%%%%%%%%%%%%%%%%%%%%%%%%%%%%%%%%%%%%
% vim: expandtab tabstop=4 shiftwidth=4 fdm=marker


\section{カーネルの有限近似}
\titlebar
% TeX source
%
% Author: Tetsuya Ishikawa <tiskw111@gmail.com>
% Date  : March 01, 2020
%%%%%%%%%%%%%%%%%%%%%%%%%%%%%%%%%%%%%%%%%%%%%%%%%%%%% SOURCE START %%%%%%%%%%%%%%%%%%%%%%%%%%%%%%%%%%%%%%%%%%%%%%%%%%%%%

\subsection{実確率変数の期待値}

突然ではあるが，確率変数の期待値について簡単に説明をさせて頂きたい．
まずは簡単のため，集合$\{z_1, z_2, \ldots, z_I \}$上のいずれかの値を取る離散確率変数$Z$を考える．
この確率変数$Z$の期待値は
\begin{equation}
    \mathrm{E}[X] = \sum_{i = 1}^{I} p_i z_i,
\end{equation}
で与えられる．ただし$p_i$は$Z$が値$z_i$を取る確率である．
連続確率変数でも同じことが言える．確率変数$X$をとある確率空間上で可測な実数値確率変数とし，
$X$が実数値$x$を取る確率密度関数を$p(x)$とおく．このとき確率変数$X$の期待値は
\begin{equation}
    \mathrm{E}[X] = \int_{\mathbb{R}} p(x) x \, \mathrm{d}x,
    \label{eqn:random_variable_expectation}
\end{equation}
で与えることができる．まさに離散確率変数の期待値を連続化したような式になっている．
ちなみに確率変数$X$が$M$次元のベクトル確率変数であった場合は
\begin{equation}
    \mathrm{E}[X] = \idotsint_{\mathbb{R}} p(\bs{x}) \bs{x}
    \, \mathrm{d}x_1 \cdots \mathrm{d}x_M,
    \label{eqn:random_variable_expectation_multi_dim}
\end{equation}
となる．

さらに，式(\ref{eqn:random_variable_expectation})の右辺を解析的に解くことが困難であったとしよう．
このとき我々には何が出来るであろうか．そう，期待値なのであるから，実際に確率変数$X$をいくつか生成してみて，
その平均値をもって期待値の近似値とすることができる．これはいわゆるMonte Carlo法であり，数式で表現すると
\begin{equation}
    \mathrm{E}[X] \approx \frac{1}{L} \sum_{\ell = 1}^{L} X_{\ell},
\end{equation}
となる．ただし$X_{\ell}$は確率変数$X$のサンプリングである．

何でいきなりこんな話をしたかって？
マァマァ，黙って式(\ref{eqn:random_variable_expectation}), (\ref{eqn:random_variable_expectation_multi_dim})
を目に焼き付けておいて頂きたい．

\subsection{カーネルのFourier変換}

さて，本題のRFFに話を戻そう．
以下ではカーネル関数$k(\bs{x}, \bs{y})$を$k(\bs{x} - \bs{y})$と表現できるものだけに限定しよう．
多項式カーネルやFisherカーネルなど多くの有名なカーネル関数が議論の対象外になってしまうものの，
式(\ref{eqn:kernel_rbf})で与えられるRBFカーネルは幸いにも$k(\bs{x} - \bs{y})$はそのように書き表すことができる．

さて，カーネル関数$k(\bs{x} - \bs{y})$をFourier変換し，その結果を$p(\bs{\omega})$と置いてみよう．すなわち
\begin{align}
    p(\bs{\omega})
    &= \mathcal{F} [k](\bs{\omega}) \notag \\
    &= \idotsint_{\mathbb{R}^M} k(\bs{z}) e^{-i\bs{\omega}\tran\bs{z}} \, \mathrm{d}z_1 \cdots \mathrm{d}z_M,
    \label{eqn:fourier_of_kernel}
\end{align}
である．これを逆Fourier変換を用いて$k(\bs{x} - \bs{y}) = \cdots$の形に書き直すと
\begin{align}
    k(\bs{x} - \bs{y})
    &= \mathcal{F}^{-1} [p](\bs{x} - \bs{y}) \notag \\
    &= \idotsint_{\mathbb{R}^M} p(\bs{\omega}) e^{i\bs{\omega}\tran(\bs{x} - \bs{y})} \, \mathrm{d}\omega_1 \cdots \mathrm{d}\omega_M,
    \label{eqn:inv_fourier_of_kernel}
\end{align}
が成り立つ．これは逆Fourier変換の定理から明らかであろう．

さて，式(\ref{eqn:inv_fourier_of_kernel})を見て思うことはないだろうか．
そう，期待値である．つまり$p$を確率密度関数とみなせば
\begin{equation}
    k(\bs{x} - \bs{y}) = \mathrm{E}_{\bs{\omega} \sim p} \left[ e^{i\bs{\omega}\tran(\bs{x} - \bs{y})} \right],
    \label{eqn:rff_expectation}
\end{equation}
と表すことができる．
ただし，上式の期待値は$\bs{\omega}$を確率変数とみなし，
さらにその確率分布は確率密度関数$p(\bs{\omega})$に従うとみなしていることに注意せよ．
期待値記号$\mathrm{E}$の添字に${\bs{\omega} \sim p}$を書いたのは，そのことを明記するためである．

先に指摘した通り，我々は期待値計算をMonte Carlo法で近似できることを知っている．
確率密度関数$p(\bs{\omega})$にしたがって確率変数$\bs{\omega}$をサンプリングした値を
$\{\bs{\omega}_\ell\}_{\ell=1}^{L}$とすると，
\begin{align}
    k(\bs{x} - \bs{y})
    &= \mathrm{E}_{\bs{\omega} \sim p} \left[ e^{i\bs{\omega}\tran(\bs{x} - \bs{y})} \right]
    \notag \\
    &\approx \frac{1}{L} \sum_{\ell = 1}^{L} e^{i\bs{\omega}_\ell\tran(\bs{x} - \bs{y})},
    \label{eqn:rff_approx}
\end{align}
が成り立つのである．
これにより，カーネル関数を，Fourier変換と乱数を用いて近似できたことになる．これがRFFの本質である．

さて，興奮に任せて大切なポイントをひとつ取り落してしまったことは読者もお気付きであろう．
式(\ref{eqn:rff_expectation})が成立するためには，関数$p(\bs{\omega})$が確率密度関数と同格でなければならない．
すなわち，可積分かつ非負の値を取る関数でなければならない．
しかし驚くべきことに，カーネル関数が正定値関数である限り，そのFourier変換は可積分かつ非負であることが証明できる．
これは非常に大事な事実であるから，定理としてまとめておく．ただし，議論の流れを妨げないよう証明は本文書の最後の節で行う．

\begin{theorem}[平行移動不変な正定値カーネルのFourier変換]
    正定値カーネル関数$k(\bs{x}, \bs{y})$が平行移動不変，すなわち$k(\bs{x}, \bs{y}) = k_{\Delta}(\bs{x} - \bs{y})$
    をみたす関数$k_{\Delta}$が存在したとする．このとき，カーネル関数$k_{\Delta}(\bs{x} - \bs{y})$のFourier変換は
    $L_1$ノルムの意味で可積分であり，かつ非負関数である．
\end{theorem}

さて，最後に近似したカーネル関数をどのように使うかについて説明しよう．
一般に，カーネル関数は正定値関数であるから，式(\ref{eqn:rff_approx})の最右辺の虚部は不要である．したがって
\begin{align}
    k(\bs{x} - \bs{y})
    &\approx \frac{1}{L} \sum_{\ell = 1}^{L} e^{i\bs{\omega}_\ell\tran(\bs{x} - \bs{y})}
    \label{eqn:rff_feature_expand_1} \\
    &= \frac{1}{L} \sum_{\ell = 1}^{L} \cos (\bs{\omega}_\ell\tran(\bs{x} - \bs{y}))
    \label{eqn:rff_feature_expand_2} \\
    &= \frac{1}{L} \sum_{\ell = 1}^{L} \bigl[
        \cos (\bs{\omega}_\ell\tran\bs{x}) \cos (\bs{\omega}_\ell\tran\bs{y}) \notag \\
    &\hspace{60pt} + \sin (\bs{\omega}_\ell\tran\bs{x}) \sin (\bs{\omega}_\ell\tran\bs{y}) \bigr],
    \label{eqn:rff_feature_expand_3}
\end{align}
を得る．ただし式(\ref{eqn:rff_feature_expand_1})から式(\ref{eqn:rff_feature_expand_2})へはEulerの公式を，
式(\ref{eqn:rff_feature_expand_2})から式(\ref{eqn:rff_feature_expand_3})へはかの懐しき三角関数の加法定理を用いた．
さてここでベクトル
\begin{equation}
    \bs{u}(\bs{x}) \triangleq \frac{1}{\sqrt{L}}
    \begin{pmatrix}
        \cos (\bs{\omega}_1\tran\bs{x}) \\
        \vdots \\
        \cos (\bs{\omega}_L\tran\bs{x}) \\[5pt]
        \sin (\bs{\omega}_1\tran\bs{x}) \\
        \vdots \\
        \sin (\bs{\omega}_L\tran\bs{x})
    \end{pmatrix}
    =
    \begin{pmatrix}
        \cos (\bs{W}\bs{x}) \\[2pt]
        \sin (\bs{W}\bs{x})
    \end{pmatrix}
    \label{eqn:rff_transform}
\end{equation}
を導入しよう．
ただし行列$\bs{W} = (\bs{\omega}_1, \ldots, \bs{\omega}_L)\tran$であり，
上式の最右辺ではベクトルの余弦関数を$\cos \bs{x} = (\cos x_1, \ldots, \cos x_L)\tran$と定義して使用している．
$\sin$についても同様である．
すると式(\ref{eqn:rff_feature_expand_3})の右辺は
\begin{equation}
    \bs{u}(\bs{x})\tran \bs{u}(\bs{y}),
\end{equation}
と非常にシンプルに書き改めることができる．すなわち
\begin{equation}
    k(\bs{x} - \bs{y}) \approx \bs{u}(\bs{x})\tran \bs{u}(\bs{y}),
    \label{eqn:rff_final}
\end{equation}
と近似することができたのである．

そもそもカーネル関数は，特徴量の内積であったことを思い出してほしい．
式(\ref{eqn:rff_final})は，カーネル関数を近似することと，
特徴量として$\phi_{\bs{x}}$ではなく$\bs{u}(\bs{x})$を使うことが等価だと言っているのである．
つまり図\ref{fig:rff_and_feature_extraction}のような構図が成り立つ．

\begin{figure}[t]
    \centerline{\includegraphics[width=200pt]{figures/rff_and_feature_extraction.pdf}}
    \caption{RFFと特徴抽出}
    \label{fig:rff_and_feature_extraction}
\end{figure}

今，特徴量はもはや関数ではなく，ただの$2L$次元の実ベクトルである．
したがって，回帰や分類などのタスクを実行するのにミョーチクリンなトリックはもはや必要ない．
有限次元の線形回帰なり線形サポートベクターマシンなりを適用すれば良いのである．
線形サポートベクターマシンの学習および推論は，カーネルサポートベクターマシンと比較すると桁違いに速い．
しかも特徴量は，高々有限次元の近似とはいえ，カーネル関数の柔軟性を残している．
このようにしてカーネル法はその柔軟性を残しつつ大規模データに適用しうるだけの高速性を手にすることができるのである．

\subsection{RBFカーネルの場合}

さて，実際にRFFを運用するためには，カーネル関数を何らかの関数に定め，
そのカーネル関数に対して確率密度関数$p(\bs{\omega})$を算出しなければならない．
RFFはカーネル関数が$k(\bs{x} - \bs{y})$と表現できるものに限られるから，ここではRBFカーネルの場合について考察しよう．

とはいえ，確率密度関数$p(\bs{\omega})$の導出は式(\ref{eqn:fourier_of_kernel})にしたがってFourier変換を
計算するだけである．直接計算だけであるので証明は省略するが，
\begin{equation}
    p(\bs{\omega}) = \frac{1}{\sqrt{2 \pi \sigma^2}^{M}} \exp \left( - \frac{\| \bs{\omega} \|^2}{2 \sigma^2} \right),
\end{equation}
となる．つまり確率変数$\bs{\omega}$は正規分布するのである．
したがって確率変数$\bs{\omega}$は正規分布からサンプリングすれば良いということになる．

以上をまとめると，RBFベースのRFFは次の手順で実施すれば良い．
\begin{enumerate}
    \item RFFの自由度$L \in \mathbb{Z}^{+}$を適当な数に設定する．これはカーネル関数を近似する際のサンプリング数である．
    \item 行列の各要素が正規分布にしたがうランダム行列$\bs{W} \in \mathbb{R}^{L \times M}$を生成する．
    \item 各学習データ$\bs{x}_i$に式(\ref{eqn:rff_transform})の変換を施し，それを改めて学習データセットとみなす．
    \item 変換後のデータセットに対し，線形回帰や線形サポートベクターマシンなどのアルゴリズム$f$を適用する．
    \item 推論時には，推論データ$\bs{x}$に対してもデータ変換式(\ref{eqn:rff_transform})を適用して$\bs{u}(\bs{x})$を算出し，
          その後に推論$f(\bs{u}(\bs{x}))$を行う．
\end{enumerate}

\subsection{数値実験：RBFカーネルのRFF近似}

本節では，RFFを用いてRBFカーネルが近似される様子を実際に確認してみよう．
まずは近似前のRBFカーネルを
\begin{equation}
    k(x_1, x_2) \triangleq \exp \left(- \frac{\| x_1 - x_2 \|^2}{2} \right),
\end{equation}
とおくとしよう．このとき対応するRFFの確率密度関数は
\begin{equation}
    p(\omega) = \frac{1}{\sqrt{2 \pi}} \exp \left(- \frac{\| \omega \|^2}{2} \right),
\end{equation}
となるので，RFF近似したRBFカーネルは
\begin{equation}
    k_{\textrm{rff}}(x_1, x_2) \triangleq \bs{u}(x_1)\tran \bs{u}(x_2),
    \hspace{10pt}
    \omega_1, \ldots, \omega_L \sim p(\omega),
\end{equation}
となる．これを試みに$L = 10, 100, 1000$として描画してみたのが図\ref{fig:rff_approximation}である．
RFFの次元を上げるにつれてRBFカーネルに近づいていくことが分かるであろう．

\begin{figure*}[t]
    \centerline{\includegraphics[width=540pt]{figures/kernel_approximation.pdf}}
    \caption{RBFカーネルのRFF近似}
    \label{fig:rff_approximation}
\end{figure*}

\subsection{深層学習との関係}

最後に，RFFを用いたサポートベクターマシンと深層学習との関係性について述べておく．
RFFは入力データ$\bs{x}$の特徴量を式(\ref{eqn:rff_transform})で定める手法であった．
これを線形サポートベクターマシンと組み合わせてみよう．

問題設定としては，教師付き学習データを用いて$C$クラスのクラス分類器を作成することを目標としよう．
学習データセットは$\mathcal{D} = \{ (\bs{x}_n, y_n) \}_{n=1}^{N}$とする．
ただし$\bs{x}_n \in \mathbb{R}^M$, $y_n \in \{1, 2, \ldots, C\}$であり，
$\bs{x}_n$が学習データに，$y_n$が教師データに対応する．

さて，線形サポートベクターマシンの判別式は線形関数であるから，
入力データ$\bs{x}$に対する各クラスの対数尤度は
\begin{equation}
    f_c(\bs{x}) = \bs{a}_c\tran \bs{x} + b_c,
    \label{eqn:rff_svm_class_wise}
\end{equation}
で表現される．ベクトル$\bs{a}_c$およびスカラー$b_c$は学習パラメータであり，各クラスの対数尤度を上手く表現できるように学習される．

さて，少しだけ表現を変更しよう．まず，すべてのクラスをまとめて記述するために，
$\bs{A} = (\bs{a}_1, \ldots, \bs{a}_C)\tran$, $\bs{b} = (b_1, \ldots, b_C)$とおく．
さらに，RFFの適用を前提とすれば，入力データは$\bs{x}$ではなく$\bs{u}(\bs{x})$に置き換えなければならない．
すると式(\ref{eqn:rff_svm_class_wise})は
\begin{equation}
    \bs{f}(\bs{x}) = \bs{A}
    \begin{pmatrix}
        \cos(\bs{Wx}) \\ \sin(\bs{Wx})
    \end{pmatrix}
    + \bs{b},
\end{equation}
となる．

これを深層学習の立場から見ると，ちょっと特殊な活性化関数を用いた，
2層の全結合層を持つニューラルネットワークとみなすことができる．
図\ref{fig:rff_svc_as_dnn}を参照されたい．

\begin{figure*}[t]
    \centerline{\includegraphics[width=350pt]{figures/rff_svc_as_dnn.pdf}}
    \caption{RFFとニューラルネットワーク}
    \label{fig:rff_svc_as_dnn}
\end{figure*}

もちろん，学習の仕方はサポートベクターマシンと深層学習で全く異なることに注意されたい．
サポートベクターマシンは区分的線形な目的関数を制約付き2次計画問題に帰着させているのに対し，
深層学習ではクロスエントロピーを確率的勾配法で最小化させている．

%%%%%%%%%%%%%%%%%%%%%%%%%%%%%%%%%%%%%%%%%%%%%%%%%%%%% SOURCE FINISH %%%%%%%%%%%%%%%%%%%%%%%%%%%%%%%%%%%%%%%%%%%%%%%%%%%%
% vim: expandtab tabstop=4 shiftwidth=4 fdm=marker


\section{数値実験}
\titlebar
% TeX source
%
% Author: Tetsuya Ishikawa <tiskw111@gmail.com>
% Date  : March 01, 2020
%%%%%%%%%%%%%%%%%%%%%%%%%%%%%%%%%%%% SOURCE START %%%%%%%%%%%%%%%%%%%%%%%%%%%%%%%%%%%

GitHubにRFFの実装およびそれを用いたいくつかの計算例を掲載してある．以下のレポジトリを参照されたい．

\verb|https://github.com/tiskw/|
\mbox{}\hspace*{20pt}\verb|Random-Fourier-Features/tree/master/|
\mbox{}\hspace*{30pt}\verb|examples/rff_svc_for_mnistsubsection|

%%%%%%%%%%%%%%%%%%%%%%%%%%%%%%%%%%%% SOURCE FINISH %%%%%%%%%%%%%%%%%%%%%%%%%%%%%%%%%%
% vim: expandtab shiftwidth=4 tabstop=4 filetype=tex


\section{定理の証明}
\titlebar
% TeX source
%
% Author: Tetsuya Ishikawa <tiskw111@gmail.com>
% Date  : October 03, 2020
%%%%%%%%%%%%%%%%%%%%%%%%%%%%%%%%%%%%%%%%%%%%%%%%%%%%% SOURCE START %%%%%%%%%%%%%%%%%%%%%%%%%%%%%%%%%%%%%%%%%%%%%%%%%%%%%

本節では，本文書で証明せずに残してあった定理を証明する．
前節までは，原論文に合わせるために角周波数ベースのFourier変換を用いて議論をしていたが，
以下では数学的シンプルさを優先して周波数ベースのFourier変換を用いることとする．

\subsection*{Fourier変換とParsevalの等式}

まずはFourier変換について簡単におさらいしておこう．本節ではFourier変換の定義として
\begin{equation}
    \mathcal{F}[f](\xi) \triangleq \int_{-\infty}^{\infty}
    f(x) e^{-2\pi i x \xi} \, \mathrm{d}x,
    \label{eqn:def_fourier_transform}
\end{equation}
を採用する．この時，逆Fourie変換は
\begin{equation}
    \mathcal{F}^{-1}[\widehat{f}](x) \triangleq \int_{-\infty}^{\infty}
    \widehat{f}(\xi) e^{2\pi i x \xi} \, \mathrm{d}\xi,
\end{equation}
となる．原論文で用いられている角周波数ベースのFourier変換とはスカラー倍の違いがあるだけで，本質的には同一である．
多変数関数$f(\bs{x}), \bs{x} \in \mathbb{R}^M$のFourier変換および逆Fourie変換は
\begin{align}
    \mathcal{F}[f](\bs{\xi}) &\triangleq \int_{-\infty}^{\infty}
    f(x) e^{-2\pi i \bs{x}\tran \bs{\xi}} \, \mathrm{d}\bs{x}, \\
    \mathcal{F}^{-1}[\widehat{f}](\bs{x}) &\triangleq \int_{-\infty}^{\infty}
    \widehat{f}(\bs{\xi}) e^{2\pi i \bs{x}\tran \bs{\xi}} \, \mathrm{d}\bs{\xi},
\end{align}
となる．ただし記号$\mathrm{d}\bs{x}$は$\mathrm{d}x_1 \mathrm{d}x_2 \cdots \mathrm{d}x_M$
の略記であることに注意せよ．

\begin{displayquote}\footnotesize\textsf{NOTE:}
    Lebesgue測度論の慣習に従えば，この$\mathrm{d}\bs{x}$は本来
    $\mu(\mathrm{d}\bs{x})$あるいは$\mathrm{d}\mu(\bs{x})$と表記されるべきものである．
    ただし$\mu$はLebesgue測度である．だが表記の冗長さと，測度論に立ち入り結論に辿り着くまでの時間が
    長くなることで読者の興味が薄れることを恐れ，ここでは上記の略記を用いた．
    ゴマカシだという謗りは受けて然るべきである．より良いアプローチがあればぜひご連絡頂きたい．
\end{displayquote}

オイオイ，式(\ref{eqn:def_fourier_transform})が定義できる関数$f$なんてかなり少なくないかと思ったそこの読者，
非常に良い眼をお持ちである．式(\ref{eqn:def_fourier_transform})の収束性については本文書末尾の付録で簡単に言及するが，
応用上必要な関数に対しては，様々なテクニックを駆使することで収束させることができるので安心して頂きたい．

次にParsevalの等式 (\textit{Parseval's identity}) について見ていこう．
Parsevalの等式と言うと公式$\|f\|^2 = \|\widehat{f}\|^2$を指すことが多いが，本文書ではより一般的な形の定理を紹介する．

\begin{theorem}[Parsevalの等式]
    任意のFourier変換可能な関数$f$, $g$およびそのFourier変換$\widehat{f}$, $\widehat{g}$に対して
    $\langle f, g \rangle = \langle \widehat{f}, \widehat{g} \rangle$が成り立つ．
\end{theorem}

ここではParsevalの等式の証明は省略する．
証明が気になる読者は（そうあるべきである！）本文書末尾の付録を参照されたい．
聡明な読者には無用の言であろうが，$g = f$とすれば$\|f\|^2 = \|\widehat{f}\|^2$を得る．

\subsection*{定理の証明}

定理を再掲し証明しよう．

\begin{theorem}[平行移動不変な正定値カーネルのFourier変換]
    正定値カーネル関数$k(\bs{x}, \bs{y})$が平行移動不変，すなわち$k(\bs{x}, \bs{y}) = k_{\Delta}(\bs{x} - \bs{y})$
    をみたす関数$k_{\Delta}$が存在したとする．このとき，カーネル関数$k_{\Delta}(\bs{x} - \bs{y})$のFourier変換は
    $L_1$ノルムの意味で可積分であり，かつ非負関数である．
\end{theorem}
\begin{proof}
    以下，$\bs{x}$, $\bs{y}$の次元を$M$とし，$\mathcal{X} = \mathbb{R}^M$とおく．すなわち$\bs{x}, \bs{y} \in \mathcal{X}$である．
    カーネル関数$k$の対称性より
    \begin{equation}
        k_{\Delta}(\bs{x}) = k(\bs{x}/2, -\bs{x}/2) = k(-\bs{x}/2, \bs{x}/2) = k_{\Delta}(-\bs{x})
    \end{equation}
    が成り立つ．これを用いれば$k_{\Delta}$のFourier変換は実数値関数であることがわかる．
    なぜならば$k_{\Delta}$のFourier変換の複素共役を計算してみると
    \begin{align}
        \mathcal{F}[k_{\Delta}]^*(\bs{\xi})
        &= \int_{\mathcal{X}} k_{\Delta}(\bs{x}) e^{2\pi i \bs{x}\tran\bs{\xi}} \, \mathrm{d}\bs{x} \\
    \intertext{ここで$\bs{x} = - \bs{z}$と置換すれば}
        &= \int_{\mathcal{X}} k_{\Delta}(-\bs{z}) e^{-2\pi i \bs{z}\tran\bs{\xi}} \, \mathrm{d}\bs{z} \\
        &= \int_{\mathcal{X}} k_{\Delta}(\bs{z}) e^{-2\pi i \bs{z}\tran\bs{\xi}} \, \mathrm{d}\bs{z}
        = \mathcal{F}[k_{\Delta}](\bs{\xi}),
    \end{align}
    となり，複素共役をとる前のFourier変換と一致する．したがって$\mathcal{F}[k_{\Delta}](\bs{\xi})$は少なくとも実数値関数である．
    次に，カーネル関数$k$が正定値関数であることから，任意の実数値関数$y$について
    \begin{equation}
        \iint_{\mathcal{X} \times \mathcal{X}} y(\bs{x}_1) k(\bs{x}_1, \bs{x}_2) y(\bs{x}_2)
        \, \mathrm{d}\bs{x}_1 \mathrm{d}\bs{x}_2 \geq 0,
    \end{equation}
    が成り立つ．上式の左辺を変形すると
    \begin{align}
        & \iint_{\mathcal{X} \times \mathcal{X}} y(\bs{x}_1) k(\bs{x}_1, \bs{x}_2) y(\bs{x}_2)
        \, \mathrm{d}\bs{x}_1 \mathrm{d}\bs{x}_2 \\
        =& \int_{\mathcal{X}} y(\bs{x}_1) \langle k(\bs{x}_1, \cdot), y \rangle \, \mathrm{d}\bs{x}_1 \\
    \intertext{ここで内積$\langle k(\bs{x}_1, \cdot), y \rangle$に対してパーセバルの不等式を用いる．
    新たに$\widehat{k}_{\bs{x}_1}(\bs{\xi}_2) = \mathcal{F}[k(\bs{x}_1, \cdot)](\bs{\xi}_2)$とおくと}
        =& \int_{\mathcal{X}} y(\bs{x}_1) \langle \widehat{k}_{\bs{x}_1}, \widehat{y} \rangle \, \mathrm{d}\bs{x}_1 \\
        =& \iint_{\mathcal{X} \times \mathcal{X}} y(\bs{x}_1) \widehat{k}_{\bs{x}_1}(\bs{\xi}_2) \widehat{y}^*(\bs{\xi}_2) \, \mathrm{d}\bs{x}_1 \mathrm{d}\bs{\xi}_2 \\
    \intertext{ここで見易さのために関数$\widehat{k}_{\bs{x}_1}(\bs{\xi}_2) = l_{\bs{\xi}_2}(\bs{x}_1)$とおくと}
        =& \iint_{\mathcal{X} \times \mathcal{X}} y(\bs{x}_1) l_{\bs{\xi}_2}(\bs{x}_1) \widehat{y}^*(\bs{\xi}_2) \, \mathrm{d}\bs{x}_1 \mathrm{d}\bs{\xi}_2 \\
        =& \int_{\mathcal{X}} \langle y, l_{\bs{\xi}_2} \rangle \widehat{y}^*(\bs{\xi}_2) \, \mathrm{d}\bs{\xi}_2 \\
    \intertext{ここで再び$\langle y, l_{\bs{\xi}_2} \rangle$に対してパーセバルの不等式を用いる．
        新たに$\widehat{l}_{\bs{\xi}_2}(\bs{\xi}_1) = \mathcal{F}[l_{\bs{\xi}_2}](\bs{\xi}_1)$とおくと}
        =& \int_{\mathcal{X}} \langle \widehat{y}, \widehat{l}_{\bs{\xi}_2} \rangle \widehat{y}^*(\bs{\xi}_2) \, \mathrm{d}\bs{\xi}_2 \\
        =& \iint_{\mathcal{X} \times \mathcal{X}} \widehat{y}(\bs{\xi}_1) \widehat{l}_{\bs{\xi}_2}(\bs{\xi}_1) \widehat{y}^*(\bs{\xi}_2) \, \mathrm{d}\bs{\xi}_1 \mathrm{d}\bs{\xi}_2,
        \label{eqn:proof_tag1}
    \end{align}
    となる．ここで$\widehat{l}_{\bs{\xi}_2}(\bs{\xi}_1)$を真面目に計算してみると
    \begin{align}
        \widehat{l}_{\bs{\xi}_2}(\bs{\xi}_1)
        &= \iint_{\mathcal{X} \times \mathcal{X}} k(\bs{x}_1, \bs{x}_2) e^{- 2\pi i (\bs{x}_1\tran\bs{\xi}_1 - \bs{x}_2\tran\bs{\xi}_2)}
        \, \mathrm{d}\bs{x}_1 \mathrm{d}\bs{x}_2 \\
        &= \iint_{\mathcal{X} \times \mathcal{X}} k_{\Delta}(\bs{x}_1 - \bs{x}_2) e^{- 2\pi i (\bs{x}_1\tran\bs{\xi}_1 - \bs{x}_2\tran\bs{\xi}_2)}
        \, \mathrm{d}\bs{x}_1 \mathrm{d}\bs{x}_2 \notag \\
        &= \delta(\bs{\xi}_1 - \bs{\xi}_2) \widehat{k}(\bs{\xi}_1),
    \end{align}
    となる．これを式(\ref{eqn:proof_tag1})に戻すと
    \begin{align}
        =& \iint_{\mathcal{X} \times \mathcal{X}}  \widehat{y}(\bs{\xi}_1)
        \delta(\bs{\xi}_1 - \bs{\xi}_2) \widehat{k}(\bs{\xi}_1) \widehat{y}^*(\bs{\xi}_2) \, \mathrm{d}\bs{\xi}_1 \mathrm{d}\bs{\xi}_2 \\
        =& \int_{\mathcal{X}} \widehat{y}(\bs{\xi}) \widehat{k}(\bs{\xi}) \widehat{y}^*(\bs{\xi}) \, \mathrm{d}\bs{\xi} \\
        =& \int_{\mathcal{X}} \| \widehat{y}(\bs{\xi}) \|^2 \, \widehat{k}(\bs{\xi}) \, \mathrm{d}\bs{\xi} \geq 0,
    \end{align}
    となる．これが任意の関数$y$について成立するから，$\widehat{y}(\bs{\xi})$が定数関数の場合も成立しなければならない．
    したがって
    \begin{equation}
        \int_{\mathcal{X}} \widehat{k}(\bs{\xi}) \, \mathrm{d}\bs{\xi} \geq 0,
    \end{equation}
    が成り立つ必要がある．したがって$\widehat{k}$は$L_1$ノルムの意味で可積分である．同様に関数$\widehat{y}(\bs{\xi})$が
    限りなくデルタ関数に近い形状の関数に対しても式(\ref{eqn:proof_tag1})が成立しなければならないから，任意の$\bs{\xi}$に対して
    $\widehat{k}(\bs{\xi}) \geq 0$が成り立つ必要がある．したがって$\widehat{k}(\bs{\xi}) \geq 0$は非負関数である．
\end{proof}

%%%%%%%%%%%%%%%%%%%%%%%%%%%%%%%%%%%%%%%%%%%%%%%%%%%%% SOURCE FINISH %%%%%%%%%%%%%%%%%%%%%%%%%%%%%%%%%%%%%%%%%%%%%%%%%%%%
% vim: expandtab tabstop=4 shiftwidth=4 fdm=marker


\section*{付録：Fourier変換}
\titlebar
% TeX source
%
% Author: Tetsuya Ishikawa <tiskw111@gmail.com>
% Date  : October 03, 2020
%%%%%%%%%%%%%%%%%%%%%%%%%%%%%%%%%%%%%%%%%%%%%%%%%%%%% SOURCE START %%%%%%%%%%%%%%%%%%%%%%%%%%%%%%%%%%%%%%%%%%%%%%%%%%%%%

本節では，定理証明の際に省略したFourier変換およびParsevalの等式の解説を行う．
本節でも数学的シンプルさを優先して周波数ベースのFourier変換を用いることとする．

\subsection*{Fourier変換に関する補足}

Fourier変換の定義としては本文の式(\ref{eqn:def_fourier_transform})で十分だが，
より良い表記，より深い理解のために，Fourier変換を関数空間の内積を用いて表現してみよう．
まず関数$\eta_{\xi}(x)$を
\begin{equation}
    \eta_{\xi}(x) \triangleq e^{2\pi i x \xi},
\end{equation}
と定義する．次に関数の内積$\langle \cdot, \cdot \rangle$を最も標準的な内積演算
\begin{equation}
    \langle f, g \rangle \triangleq \int_{-\infty}^{\infty} f(x) g^*(x) \, \mathrm{d}x,
\end{equation}
で定めるとしよう．ただし$g^*(x)$は関数$g(x)$の複素共役である．
このとき式(\ref{eqn:def_fourier_transform})で表現でされるFourier変換は
\begin{equation}
    \mathcal{F}[f](\xi) = \langle f, \eta_{\xi} \rangle,
\end{equation}
と表すことができる．
このとき関数群$\{ \eta_{\xi}(x) \,|\, \xi \in \mathbb{R} \}$は正規直交基底と非常に似た性質を持つ．すなわち
\begin{equation}
    \langle \eta_{\xi_1}, \eta_{\xi_2} \rangle = \delta(\xi_1 - \xi_2),
    \label{eqn:fourier_base}
\end{equation}
が成り立つ．ただし$\delta$はDiracのデルタ関数 (\textit{Dirac delta function}) である．
つまり式(\ref{eqn:fourier_base})は$\xi_1 = \xi_2$のときのみ値を持つ．
これは通常のベクトル空間の基底$\{ \bs{e}_1, \bs{e}_2, \ldots \}$が持つ性質
\begin{equation}
    \bs{e}_i\tran \bs{e}_j =
    \begin{cases}
        1, & i = j, \\
        0, & i \neq j,
    \end{cases}
\end{equation}
と非常に良く似ている．
したがってFourier変換は正規直交基底との内積，つまるところただの基底交換のようなもの，と見ることもできるのである．

\begin{displayquote}\footnotesize\textsf{NOTE:}
    Fourier変換の定義を初めて見た若かりし頃，諸君らも
    「何でFourier変換が逆回転$e^{-2\pi i x \xi}$で逆Fourier変換が順回転$e^{2\pi i x \xi}$なんだよ，フツーは逆だろ！」
    とキレかかった経験があるのではないだろうか．上述の内積によるFourier変換の理解はこの疑問を解消してくれる．
    そう，基底$\eta_{\xi}(x)$はちゃんと順回転$e^{2\pi i x \xi}$で定義されていたのだ．
    しかし内積を計算する際に共役をとる必要があり，そこで回転方向が反転したのである．少しはスッキリしたであろうか．
\end{displayquote}

\subsection*{Fourier変換の収束性について}

Parsevalの等式の証明に移る前に，一部の読者が気になっているであろう
Fourier変換(\ref{eqn:def_fourier_transform})の収束性について簡単に述べておこう．
つまり，式(\ref{eqn:def_fourier_transform})が定義できるのはどんな関数$f$か？という問題である．
結論から言えば，式(\ref{eqn:def_fourier_transform})が収束する条件は
\begin{equation}
    \int_{-\infty}^{\infty} |f(x)| \, \mathrm{d}x < \infty,
\end{equation}
である．ハッキリ言って，上式を満たす関数は応用上の観点から見ても非常に少ない．
例えば三角関数や指数関数ですら上式の制約は満たさない．
もちろん世の中には病的な関数が沢山あるゆえ，真に全ての関数に対してFourier変換を定義せよというのは困難であるのだが，
とはいえ，少なくとも応用上必要な関数に対してはFourier変換ができないと困ってしまう．
この問題に対するアプローチは様々あるが，著者の知る限りでは，
\begin{enumerate}
    \item[(1)] 扱う関数はすべて有限の台 (\textit{support}) を持つと仮定する方法．
               すなわち，関数が値を持つ範囲を十分に広い有限な範囲に限定してしまうという方法．
               関数の台が有限であれば，応用上必要な関数はほぼすべて積分可能となる．
    \item[(2)] 両側Laplace変換の特殊な場合 ($s = i\omega$) としてFourier変換を定義する方法である．
               実は両側Laplace変換は，解析接続を用いた巧妙な仕組みにより，非常に幅広い関数に対して定義可能である．
\end{enumerate}
の2つのアプローチがある．初学者は(1)の理解で良いであろう．
世の中の多くの人の理解は以下の(1)の方法で理解していると思われる．
ちなみに著者は理論的にも便利な(2)で理解している．興味があれば調べてみると良い．

\subsection*{Parsevalの等式}

次にParsevalの等式 (\textit{Parseval's identity}) について見ていこう．
一般にParsevalの等式と言うと$\|f\|^2 = \|\widehat{f}\|^2$を指すことも多いが，本文書ではより一般的な形の定理を証明する．
Fourier変換が正規直交基底を用いた基底の交換であることが分かれば，ほぼ自明な定理である．

\begin{theorem}[Parsevalの等式]
    任意のFourier変換可能な関数$f$, $g$およびそのFourier変換$\widehat{f}$, $\widehat{g}$に対して
    $\langle f, g \rangle = \langle \widehat{f}, \widehat{g} \rangle$が成り立つ．
\end{theorem}
\begin{proof}
    証明は直接計算による．
    \begin{align*}
        \langle \widehat{f}, \widehat{g} \rangle
        &= \int_{-\infty}^{-\infty} \langle f, \eta_{\xi} \rangle \langle g, 
        \eta_{\xi} \rangle \, \mathrm{d}\xi \\
        &= \iiint_{-\infty}^{-\infty}  f(x_1) g(x_2) \eta_{\xi}(x_1) \eta_{\xi}(x_2)
        \, \mathrm{d}x_1 \mathrm{d}x_2 \mathrm{d}\xi \\
        \intertext{ここで$\eta_{\xi}(x_1) \eta_{\xi}(x_2) = \eta_{x_1}(\xi) \eta_{x_2}(\xi)$に注意すれば}
        &= \iint_{-\infty}^{-\infty}  f(x_1) g(x_2)
        \langle \eta_{x_1}, \eta_{x_2} \rangle \, \mathrm{d}x_1 \mathrm{d}x_2 \\
        \intertext{関数$\eta_{x}(\xi)$の正規直交性に注意すれば
        $\langle \eta_{x_1}, \eta_{x_2} \rangle$は$\delta(x_1 - x_2)$であるから}
        &= \int_{-\infty}^{-\infty}  f(x) g(x) \, \mathrm{d}x
        = \langle f, g \rangle,
    \end{align*}
    となり，題意が成り立つ．
\end{proof}

ここまで読み進めた読者には指摘するまでもないであろうが，
上記定理の特殊な場合として$f = g$とおけば，良く知られた$\|f\|^2 = \|\widehat{f}\|^2$を得る．

%%%%%%%%%%%%%%%%%%%%%%%%%%%%%%%%%%%%%%%%%%%%%%%%%%%%% SOURCE FINISH %%%%%%%%%%%%%%%%%%%%%%%%%%%%%%%%%%%%%%%%%%%%%%%%%%%%
% vim: expandtab tabstop=4 shiftwidth=4 fdm=marker


\section*{おわりに}
\titlebar
% TeX source
%
% Author: Tetsuya Ishikawa <tiskw111@gmail.com>


TBD


% 最後に，私の数学的活動は，2017年に逝去された恩師，山下弘一郎先生や，
% 大学および大学院で私の指導教官を担当して下さった早川朋久准教授をはじめ，
% 数学で私と関わりを持ったすべての方々のおかげで成り立っています．
% そして，数学的活動の以前に，そもそも私の生は両親によって与えられています．


% vim: expandtab tabstop=4 shiftwidth=4 fdm=marker


% \pagebreak

% TeX source
%
% Author: Tetsuya Ishikawa <tiskw111@gmail.com>
% Date  : Aug  4, 2019
%%%%%%%%%%%%%%%%%%%%%%%%%%%%%%%%%%%% SOURCE START %%%%%%%%%%%%%%%%%%%%%%%%%%%%%%%%%%%


\begin{thebibliography}{9}

\bibitem{Vapnik1963}
V.~Vapnik and A.~Lerner,
``Pattern recognition using generalized portrait method'',
Automation and Remote Control, vol. 24, 1963.

\bibitem{Boser1992}
B.~Boser, I.~Guyon and V.~Vapnik,
``A training algorithm for optimal margin classifiers'',
Proc. of the fifth annual workshop on Computational learning theory, pp. 144-152, 1992

\bibitem{Akaho2008}
赤穂昭太郎，
「カーネル多変量解析」，
岩波書店，2008．

\bibitem{Jost2005}
J. Jost,
``Postmodern Analysis'',
Springer, 2005.

\bibitem{Horiuchi2005}
堀内利郎，下村勝孝，
「関数解析の基礎 -- $\infty $次元の微積分」，
内田老鶴圃，2005．

\bibitem{Bishop2010}
C. Bishop,
``Pattern Recognition and Machine Learning'',
Springer, 2010.

\bibitem{Rahimi2007}
A.~Rahimi and B.~Recht, 
``Random Features for Large-Scale Kernel Machines'',
Neural Information Processing Systems, 2007.

\bibitem{Yu2016}
F.~X.~Yu, A.~T.~Suresh, K.~Choromanski, D.~H.~Rice and S.~Kumar,
``Orthogonal Random Features'',
Neural Information Processing Systems, 2016.

\end{thebibliography}


%%%%%%%%%%%%%%%%%%%%%%%%%%%%%%%%%%%% SOURCE FINISH %%%%%%%%%%%%%%%%%%%%%%%%%%%%%%%%%%
% vim: expandtab shiftwidth=4 tabstop=4 filetype=tex


\end{document}

%%%%%%%%%%%%%%%%%%%%%%%%%%%%%%%%%%%% SOURCE FINISH %%%%%%%%%%%%%%%%%%%%%%%%%%%%%%%%%%
% vim: expandtab shiftwidth=4 tabstop=4 filetype=tex
